\documentclass[UTF8,10pt]{ctexart}
% Sprint cadence revised after Lesson 2; source retained as the canonical Gantt definition.
\usepackage[a4paper,landscape,margin=1.15cm]{geometry}
\usepackage{pgfgantt}
\usepackage{xcolor}
\usepackage{array}
\pagestyle{empty}

\definecolor{donegray}{RGB}{92,99,112}
\definecolor{planblue}{RGB}{67,105,163}
\definecolor{stretchgold}{RGB}{176,126,28}
\definecolor{docgreen}{RGB}{77,126,94}
\definecolor{gridgray}{RGB}{220,223,228}
\definecolor{currentorange}{RGB}{186,105,38}

\begin{document}
\begin{center}
{\LARGE\bfseries 数据备份软件课程项目滚动甘特图}\par
\vspace{0.15em}
{\small 以 10 次课程为检查节点；每个 Sprint 从本次课启动，到下一次课评审、回顾并收口}\par
\vspace{0.55em}
\end{center}

\begin{ganttchart}[
    hgrid style/.style={draw=gridgray},
    vgrid={*{9}{draw=gridgray}, draw=gray!55},
    x unit=1.34cm,
    y unit title=0.58cm,
    y unit chart=0.45cm,
    title height=0.95,
    title label font=\bfseries\small,
    bar height=0.48,
    group height=0.34,
    group peaks tip position=0,
    group peaks height=0.16,
    bar label font=\scriptsize,
    group label font=\bfseries\small,
    milestone label font=\scriptsize,
    milestone height=0.48,
    milestone top shift=0.15,
    bar/.append style={fill=planblue!78, draw=planblue!95},
    group/.append style={fill=planblue!18, draw=planblue!72},
    milestone/.append style={fill=donegray, draw=donegray},
]{1}{10}

\gantttitle{课程检查节点}{10} \\
\gantttitlelist{1,...,10}{1} \\

\ganttgroup[group/.append style={fill=donegray!18,draw=donegray!80}]{Sprint 迭代主线}{1}{10} \\
\ganttbar[bar/.append style={fill=currentorange!76,draw=currentorange!95}]{S1：v0.1 基础 Backup / Restore；第2次课收口}{1}{2} \\
\ganttbar{S2：v0.2 Backup Package：Archive v0.1 + Metadata v0.1 + 极简 UI}{2}{3} \\
\ganttbar{S3：v0.3 Filter + Package Hardening + 测试环境完善}{3}{4} \\
\ganttbar{S4：v0.4 程序框架 / 接口收口 + 特殊文件 + UML 同步}{4}{5} \\
\ganttbar{S5：v0.5 Archive 规范化与集成；结合打包解包课程完善}{5}{6} \\
\ganttbar{S6：v0.6 Huffman 压缩 / 解压 + Round-trip 测试}{6}{7} \\
\ganttbar{S7：v0.7 加密 / 解密 + 压缩加密流水线集成}{7}{8} \\
\ganttbar{S8：v0.8 定时备份 + inotify 实时备份 + 后台状态}{8}{9} \\
\ganttbar[bar/.append style={fill=stretchgold!72,draw=stretchgold!95}]{S9：Release Candidate；全量回归 + 增量 / 网络等 Stretch}{9}{10} \\
\ganttmilestone[milestone/.append style={fill=docgreen,draw=docgreen}]{Final Release / 答辩}{10} \\

\ganttgroup[group/.append style={fill=docgreen!18,draw=docgreen!85}]{持续工程活动}{1}{10} \\
\ganttbar[bar/.append style={fill=docgreen!68,draw=docgreen!95}]{需求 / 架构 / UML 随实现滚动更新}{1}{9} \\
\ganttbar[bar/.append style={fill=docgreen!68,draw=docgreen!95}]{自动测试 / CI / Regression：每个 Sprint 都执行}{1}{9} \\
\ganttbar[bar/.append style={fill=docgreen!68,draw=docgreen!95}]{Kanban + Sprint Backlog + Review / Retrospective}{2}{9} \\
\ganttbar[bar/.append style={fill=docgreen!68,draw=docgreen!95}]{测试报告、最终报告、演示材料持续归档}{7}{10} \\

\ganttgroup[group/.append style={fill=stretchgold!14,draw=stretchgold!78}]{主要责任线（Primary Owner）}{2}{9} \\
\ganttbar[bar/.append style={fill=stretchgold!66,draw=stretchgold!92}]{Core Algorithm \& Integration：Archive → Huffman → Encryption → Integration}{2}{9} \\
\ganttbar[bar/.append style={fill=stretchgold!52,draw=stretchgold!82}]{Architecture / Product / UI：UML、UI、交互与演示}{2}{9} \\
\ganttbar[bar/.append style={fill=stretchgold!38,draw=stretchgold!74}]{Filesystem / Engineering：Metadata、Filter、特殊文件、Automation、Quality}{2}{9}

\end{ganttchart}

\vspace{0.55em}
{\small
\begin{tabular}{@{}>{\raggedright\arraybackslash}p{0.17\linewidth} p{0.78\linewidth}@{}}
\textbf{当前节点} & 第 2 次课：完成 Sprint 1 的测试、缺陷修复、Review / Retrospective，并在同一节点启动 Sprint 2；Sprint 1 不是“第 1 次课当天完成”，而是跨第 1--2 次课的一周迭代。 \\
\textbf{每轮 DoD} & 必须形成“可运行版本 + 自动测试 / 回归结果 + 关键源码可解释 + UML / 文档同步”；测试不再集中到第 9 次课，第 9--10 次课主要用于 Release Candidate、全量回归和答辩准备。 \\
\textbf{开发策略} & 本地 120+ 分能力优先；网络、增量、多算法等属于后续 Stretch。三人按 Primary Owner 分工，但关键代码至少有 Secondary Reviewer，避免三个子项目到最后才第一次集成。 \\
\end{tabular}
}

\end{document}
